% The leantex port of vincentqb.github.io: one source, three outputs —
% index.html, llms.txt (markdown), and a print resume (PDF) — from the same
% document. Build with build-leantex.sh; compare against the Jinja build
% with compare-leantex.js. See README.md.

\documentclass{article}

\pdfmeta{
  author  = "Vincent Quenneville-Bélair",
  title   = "Vincent Quenneville-Bélair, PhD",
  subject = "Vincent Quenneville-Bélair, PhD.",
  url     = "https://vincentqb.github.io/",
  image   = "https://vincentqb.github.io/profile.jpg",
  favicon = "favicon.svg",
}

\output{
  formats = html, md,
  md = "llms.txt",
  css = own,
  stylesheet = "style.css",
}

% Inter resolves against the machine, as on the Jinja site (a Google Fonts
% load there); the fonts directory ships the Font Awesome 5 Free faces the
% icons resolve against (SIL OFL 1.1, see fonts/LICENSE-FontAwesome.txt),
% so the icons render the same on every host.
\fonts{
  dir  = "fonts",
  body = "Inter",
}

% The site's palette, corrected to hold the engine's contrast contract.
% The site's own values fail WCAG 2.2 SC 1.4.3 through the engine's
% arithmetic: primary hsl(0,85%,54%) = #ED2626 reads 4.28:1 on white and
% 4.10:1 on white-bis; Bulma's grey #69748C reads 4.49:1 on white-bis.
% Same hues, darker: primary at 45% lightness (5.40:1 / 5.17:1) and the
% grey at 44% (5.19:1 on white-bis). primary40 is Bulma's derived 40%
% shade, kept for the timeline and the button, where it passes.
\palette{
  bg        = #FFFFFF,
  primary   = #D41111,
  primary40 = #BD0F0F,
  muted     = #606A80,
  whitebis  = #FAFAFA,
}

% The menu's interaction states, declared where the menu is: hover takes
% the accent, and the transition duration is the site's own (Bulma's
% --bulma-duration, 294ms). The reduced-motion guard travels with the
% declaration by construction.
\style{nav}{
  hover  = primary,
  motion = 294ms,
}

\page{
  size = letter,
  vmargin = 0.5in,
  hmargin = 0.75in,
}

% The document title is the level-0 heading: <header><h1> in HTML, the one
% # in llms.txt, LARGE bold on the printed page.
\title{\primary{Vincent} Quenneville-Bélair, \primary{{\scshape phd}}}
\subtitle{Applied Mathematician}

\begin{document}

% ---- Menu: the site's navbar, a labeled navigation landmark ---------------
% Web-only, as on the site (its print stylesheet hides the navbar): the
% printed resume is one page and needs no in-page navigation.

\begin{ifbackend}{html}
\begin{nav}[label = Sections]
\href{#experience}{Experience}
\href{#education}{Education}
\href{#profile}{Profile}
\end{nav}
\end{ifbackend}

% ---- Hero: name, degree, title, links -------------------------------------

\maketitle

\begin{center}
  \href{mailto:vincentqb@gmail.com}{\faIcon[label = Email]{envelope}}
  \href{https://linkedin.com/in/vincentqb}{\faIcon[label = LinkedIn]{linkedin}}
  \href{https://github.com/vincentqb}{\faIcon[label = GitHub]{github}}
  \href{https://scholar.google.com/citations?user=B3XqkfMAAAAJ}{\faIcon[label = Google Scholar]{graduation-cap}}
\end{center}

% ---- Experience ------------------------------------------------------------

\section*{Experience}

\block[before = 0ex]{
  \subsection*{Machine Learning Scientist}
  \primary{Supply Chain Optimization Technologies, Amazon}\\
  \muted{\textit{2021 -- Present}}
  \begin{itemize}
    \item Build an autonomous search over optimization policies, using a Lean kernel to certify candidates and promote provably best ones.
    \item Prototype predictive model for inventory placement, accelerating pipelines by reframing decision-making into a forecasting formulation.
    \item Develop foundation forecasting model, unifying across products worldwide at multiple scales, deploying zero-shot for external partners.
    \item Design deep learning attention-based models to forecast product demand for supply chain optimization.
    \item Leverage LLM embeddings and approximate nearest-neighbor retrieval to capture cross-product demand effects, sharing demand signal across semantically similar products.
  \end{itemize}
}

\block[before = 0ex]{
  \subsection*{Machine Learning Engineer}
  \primary{PyTorch, Facebook AI}\\
  \muted{\textit{2019 -- 2021}}
  \begin{itemize}
    \item Create torchaudio library, engaging over 200 open source contributors.
    \item Develop speech-to-text and text-to-speech deep learning models.
    \item Integrate novel PyTorch optimizers collaborating with academic and industrial researchers.
  \end{itemize}
}

\block[before = 0ex]{
  \subsection*{Machine Learning Scientist}
  \primary{Supply Chain Forecasting and Marketing Measurement, Amazon}\\
  \muted{\textit{2017 -- 2019}}
  \begin{itemize}
    \item Generate long-term strategic inventory capacity forecast. Model counterfactual scenarios. Optimize supply chain expansion recommendations.
    \item Build causal impact models to estimate advertising campaign lift. Infer future returns and optimize advertising portfolio.
    \item Teach internal machine learning specialization course series on matrix factorization, dimensionality reduction, and recommender systems.
  \end{itemize}
}

\block[before = 0ex]{
  \subsection*{Chief Data Scientist}
  \primary{Vizanda}\\
  \muted{\textit{2016 -- 2017}}
  \begin{itemize}
    \item Develop web-based software to automatically extract information from structured and unstructured data on the fly. Process the information to programmatically generate relevant visualizations and intelligent insights.
  \end{itemize}
}

\block[before = 0ex]{
  \subsection*{Chu Assistant Professor of Applied Mathematics}
  \primary{Columbia University}\\
  \muted{\textit{2015 -- 2017}}
  \begin{itemize}
    \item Model gravitational wave propagation to understand black hole collisions. Create and analyze novel specialized mixed finite element methods to simulate these models numerically.
  \end{itemize}
}

% ---- Education -------------------------------------------------------------

\section*{Education}

\begin{itemize}
  \item \textbf{PhD Applied Mathematics}, \primary{University of Minnesota}\\
    \muted{\textit{NSERC Alexander Graham Bell Canada Graduate Scholarship,
    FQRNT Research Scholarship,
    University of Minnesota Doctoral Dissertation Fellowship.}}
  \item \textbf{MCS Computer Science}, \primary{University of Minnesota}
  \item \textbf{MSc Applied Mathematics}, \primary{University of Minnesota}
  \item \textbf{BSc Mathematics and Physics}, \primary{McGill University}
\end{itemize}

% ---- Profile ---------------------------------------------------------------

\section*{Profile}

\begin{center}
  \includegraphics[width=192pt, alt={Vincent Quenneville-Bélair}]{profile.jpg}

  Vincent Quenneville-Bélair, PhD, Applied Mathematician.
\end{center}

% ---- Return to top ---------------------------------------------------------
% The site's button, as declarations: a labeled landmark pinned to the
% viewport corner, revealed after 300px of scroll (the site's own JS
% threshold) — scroll-driven CSS where the platform has it, the engine's
% constant script fallback elsewhere (release Firefox), always visible
% with neither. Web-only: the printed page has no viewport.

\begin{ifbackend}{html}
\begin{nav}[label = Return to top, pin = bottom right, offset = 2em, reveal = 300px]
\href{#top}{\faIcon[label = Return to top]{arrow-up}}
\end{nav}
\end{ifbackend}

\end{document}
